\documentclass[10pt,twocolumn,a4paper]{article}

\usepackage{xetexko}
\usepackage{geometry}
\usepackage{titlesec}
\usepackage{booktabs}
\usepackage{array}
\usepackage{tabularx}
\usepackage{setspace}
\usepackage{indentfirst}
\usepackage{hyperref}
\usepackage{caption}
\usepackage{float}

\geometry{
  a4paper,
  top=25mm, bottom=25mm,
  left=18mm, right=18mm,
  columnsep=6mm
}

\setmainfont{Apple SD Gothic Neo}[
  BoldFont = Apple SD Gothic Neo Bold,
  ItalicFont = Apple SD Gothic Neo Regular
]
\setsansfont{Apple SD Gothic Neo}
\setmonofont{Courier New}
\setmainhangulfont{Apple SD Gothic Neo}[
  BoldFont = Apple SD Gothic Neo Bold
]

\titleformat{\section}[block]
  {\normalsize\bfseries\centering}
  {\Roman{section}.\,}{0.3em}{}
\titlespacing{\section}{0pt}{8pt}{4pt}

\titleformat{\subsection}[block]
  {\normalsize\bfseries}
  {\thesubsection\,}{0.3em}{}
\titlespacing{\subsection}{0pt}{6pt}{3pt}

\setlength{\parindent}{1em}
\setlength{\parskip}{2pt}

\renewcommand{\footnoterule}{\vspace{-3pt}\noindent\rule{0.4\columnwidth}{0.4pt}\vspace{3pt}}

\pagestyle{empty}

\begin{document}

\twocolumn[{%
\begin{center}
  {\Large\bfseries 실제 관광 경로 기반 AI 에이전트를 활용한}\\[4pt]
  {\Large\bfseries 제주 설화·민담 여행 코스 추천 서비스 개발}\\[6pt]
  {\normalsize --- 중간 보고서 ---}\\[10pt]

  {\normalsize 조익준$^\dagger$}\\[2pt]
  {\normalsize 경희대학교 컴퓨터공학과}\\[2pt]
  {\small\itshape e-mail: keepdev0919@gmail.com}\\[14pt]
\end{center}
\vspace{2pt}
\noindent\rule{\textwidth}{0.4pt}
\vspace{6pt}
}]

% ── 요약 ─────────────────────────────────────────────────
{\small\noindent\textbf{요 약}\quad
본 연구는 제주도 방문 관광객을 대상으로, 실제 관광객이 검증한 여행 코스에
AI가 설화·민담 콘텐츠를 매핑하여 계획·탐험·기록이 하나의 여정으로 연결되는
iOS 모바일 서비스를 구현한다. LangGraph 기반 2단계 에이전트 파이프라인이
지역·스타일·기간 입력을 받아 코스를 추천하고, 페르소나 기반 AI 설화 동행자가
현장에서 RAG 기반 설화 해설을 SSE 스트리밍으로 제공하며, 여행 후 GPT-4o가
방문 이력과 대화 맥락을 바탕으로 일지를 자동 생성한다. 데이터 파이프라인은
제주도청 공공 API에서 수집한 설화·민담을 지명 Geocoding으로 위치 정보와
연결하고 ChromaDB에 벡터화하며, Visit Jeju 실제 관광객 코스 데이터에서
중복·저품질 코스를 필터링하여 추천 후보 DB를 구성한다. 중간 보고 시점
기준으로 코스 추천·GPS 탐험·동행자 채팅·여행 일지·앱 세션 복원 기능이
모두 구현 완료되었다.}

\vspace{6pt}
{\small\noindent\textbf{키워드}\quad
LangGraph, RAG, 위치 기반 서비스, 제주 설화·민담, iOS Swift, AI 에이전트,
여행 동행자}

\vspace{8pt}

% ── I. 서론 ─────────────────────────────────────────────
\section{서 론}

제주도는 연간 1,500만 명 이상이 방문하는 국내 최대 관광지이나, 현행 관광
앱의 대부분은 맛집·숙박·교통 중심의 정보 제공에 머물고 있다. 제주도에는
설화 182건, 민담 323건 등 505건의 문화유산 콘텐츠가 공공 데이터로 개방되어
있음에도, 이를 관광 경험과 결합한 서비스는 부재한 상황이다.

최근 LLM(Large Language Model)의 발전과 함께, 특정 도메인 지식을 벡터
데이터베이스에 저장하고 검색·생성을 결합하는 RAG 기술이 주목받고 있다.
RAG는 모델의 환각(Hallucination) 문제를 억제하면서 도메인 특화 지식을
활용할 수 있어, 정확성이 요구되는 문화유산 콘텐츠 분야에 적합하다.

본 연구는 \textbf{``실제 검증된 여행 코스 위에 설화를 얹는다''} 는 컨셉으로,
제주관광공사 Visit Jeju의 실제 관광객 여행 일정 데이터를 기반으로 LangGraph
AI 에이전트가 설화·민담이 풍부한 코스를 추천하고, 여행 중에는 AI 설화 동행자
``마을 할망''이 장소에 깃든 이야기를 전달하는 iOS 앱을 구현하는 것을 목표로 한다.

% ── II. 관련 연구 ───────────────────────────────────────
\section{관련 연구}

\subsection{RAG(Retrieval-Augmented Generation)}

RAG는 Lewis et al.(2020)이 제안한 방법으로\,[1], 외부 지식 저장소에서 관련
문서를 검색(Retrieve)한 뒤 이를 문맥으로 제공하여 LLM이 더 정확한 답변을
생성(Generate)하도록 한다. 벡터 임베딩 기반의 의미론적 검색을 통해 키워드
검색의 한계를 극복하며, 특정 도메인에 특화된 지식베이스 구축에 활용된다.
본 연구는 RAG를 활용하여 ChromaDB에 저장된 제주 설화·민담을 실시간으로
검색하고, 여행 현장 컨텍스트에 맞는 해설을 생성하는 데 적용한다.

\subsection{위치 기반 문화 콘텐츠 서비스}

Pok\'{e}mon GO, Ingress 등의 위치 기반 서비스는 GPS를 활용하여 사용자가
실제 장소를 방문하도록 유도한 바 있다. 문화유산 분야에서는 Google Arts \&
Culture, 국내 문화재청의 e뮤지엄 등이 위치 기반 콘텐츠를 제공하나, 지역
설화·민담과 같은 무형 문화유산과 AI 코스 추천을 결합한 서비스는 미흡한
실정이다. 본 연구는 실제 관광객 동선 데이터와 무형 문화유산을 결합하여,
기존 서비스에서 미흡했던 설화·민담 기반 AI 코스 추천을 구현한다.

\subsection{대화형 AI 인터페이스}

ChatGPT, Claude 등의 등장 이후 사용자가 자연어로 정보를 탐색하는 대화형
인터페이스가 보편화되고 있다. RAG를 결합한 도메인 특화 챗봇은 일반 LLM의
지식 한계를 보완하면서 특정 주제에 깊이 있는 답변을 제공한다. 본 연구에서는
이를 관광 맥락에 적용하여 현장형 AI 설화 동행자 ``마을 할망''을 구현한다.

\subsection{LLM 에이전트 및 LangGraph}

Yao et al.(2023)이 제안한 ReAct(Reasoning + Acting) 패턴은\,[8] LLM이 추론
단계와 도구 사용 단계를 반복하며 복잡한 작업을 수행하도록 한다. LangGraph는\,[9]
이러한 에이전트 워크플로우를 유향 그래프(StateGraph)로 명시적으로 정의하여,
조건부 분기·루프·병렬 실행을 지원하고 에이전트의 상태를 체계적으로 관리한다.
본 연구는 LangGraph를 활용하여 코스 추천 에이전트의 분기 로직과 상태 관리를
구현한다.

% ── III. 시스템 설계 ────────────────────────────────────
\section{시스템 설계}

\subsection{전체 아키텍처}

본 시스템은 iOS 앱(Swift/SwiftUI), FastAPI 백엔드, 데이터 레이어의 3계층으로
구성된다. iOS 앱은 REST/SSE를 통해 백엔드와 통신하며, 백엔드는 LangGraph
코스 추천 에이전트, 설화 동행자 SSE 엔드포인트, GPS 핀 조회 API, TTS API로
구성된다. 데이터 레이어는 Visit Jeju 여행일정 DB(큐레이션 코스)\,[4],
ChromaDB(설화·민담 벡터 DB)\,[5], SQLite(메타데이터·GPS 좌표)로
이루어진다.

사용자 흐름은 계획·탐험·기록의 단일 여정으로 구성된다. 계획 단계에서는
취향(지역·스타일·기간·설화취향)을 입력하면 우선순위 1순위 코스가 즉시
상세 미리보기로 표시되며, 하단의 ``새로운 추천'' 버튼으로 2·3순위 코스를
순서대로 확인할 수 있다. 원하는 코스는 ``내 일정으로 담기''로 저장한다.
탐험 단계에서는
실시간 GPS 안내와 함께 장소 도착 시 AI 설화 동행자가 말을 건다.
기록 단계에서는 여행 종료 후 GPT-4o가 여행 일지를 자동으로 생성한다.
앱 재시작 시 SessionRestoreView가 진행 중인 탐험 세션을 자동 복원하여,
앱 강제 종료에도 여행 연속성을 보장한다.

\subsection{데이터 파이프라인}

\textbf{Visit Jeju 여행 데이터.}\quad 제주관광공사 Visit Jeju 공공데이터\,[4]에서
여행장소(GPS 좌표 포함)와 여행세부일정을 수집하였다. 장소명 기준 조인으로
일정별 방문 장소의 GPS 좌표를 확보한 뒤, 공항·항구 중복 도배 코스 등
저품질 코스를 필터링하여 큐레이션 코스 DB를 구성하였다.

\textbf{설화·민담 GPS 좌표 연결.}\quad 설화\,[2]는 텍스트 내 지명을 제주
행정구역 화이트리스트 기반 매칭으로 추출한 뒤 Geocoding한다\,[7].
민담\,[3]은 채록 메타데이터의 조사 장소를 파싱하여 Geocoding한다.
GPS 좌표가 확보된 항목은 OpenAI Text Embedding\,[6]을 이용해 청크로
분할하여 ChromaDB에 벡터화하고, 좌표 미확보 항목도 텍스트 검색용으로
ChromaDB에 함께 적재하였다.

\subsection{코스 추천: 2단계 API 파이프라인}

초기에는 설화 취향을 먼저 입력받아 설화에서 코스를 역설계하는 방식을
시도하였으나, 사용자가 실제로 여행을 계획할 때 지역·기간이 먼저 결정된다는
점을 반영하여 현재의 구조로 재설계하였다.

응답 지연을 최소화하기 위해 코스 추천을 두 단계로 분리하였다. 첫 번째
단계(\texttt{POST /course/list})에서는 LangGraph course\_list\_agent가 SQLite에서
후보 코스 30개를 무작위로 조회하고, GPT-4o가 사용자의 여행 스타일에 적합한 코스
3개를 선별하여 빠르게 반환한다. 지역 필터는 코스 내 장소의 50\% 이상이 해당
GPS 범위 안에 있을 때 유효한 것으로 판단한다.

두 번째 단계(\texttt{POST /course/detail})에서는 결정론적 함수 파이프라인으로
처리한다. 선택된 코스의 모든 장소에 대해 haversine 거리를 계산하여 반경
3km 이내 설화를 매핑하고, GPT-4o가 설화 취향과 장소 정보를 바탕으로 여행
내러티브 2문장을 생성한다. 이 단계는 LLM 호출이 1회만 발생하므로 예측
가능한 비용과 안정성을 갖는다.

\subsection{AI 설화 동행자 시스템}

여행 중 경험의 중심축으로 페르소나 기반 AI 설화 동행자 시스템을 구현하였다.
RAG 챗봇 로직을 독립 모듈로 분리하여 동행자 채팅 기능과 연결하였으며,
페르소나 정의(이름·말투·방언 규칙)를 별도 설정으로 관리하여 향후 다양한
동행자 캐릭터로 확장할 수 있도록 설계하였다. 동행자는 해당 장소 주변 설화를
ChromaDB에서 벡터 검색하여 맥락으로 주입받으며, FastAPI
SSE(Server-Sent Events)를 통해 토큰 단위 스트리밍으로 응답한다.

현재 제공되는 기본 페르소나는 제주 토박이 할머니 ``마을 할망''으로,
문장당 제주 방언 표현을 자연스럽게 녹이는 방식(예: ``혼저 옵서, 삼춘.
이 당에 하영 반갑수다.'')으로 현지감을 구현하였다. 멀티턴 대화 히스토리와
사용자가 저장한 코스 컨텍스트를 유지하여, 방문 장소와 설화 이야기를
연속적인 대화로 풀어낸다.

\subsection{GPS 탐험 모드 및 도착 감지}

iOS CoreLocation 기반 지오펜스로 도착을 감지한다. 도보 모드에서는 반경
100m 진입 후 30초 체류, 차량 모드에서는 반경 300m 진입 후 30초 체류 시
도착으로 판정한다. 30초 체류 조건은 잠깐 지나치는 오탐을 방지한다. 도착 판정
시 동행자 오버레이가 전체화면으로 등장하며, 사용자가 탭하면 해당 장소의 설화를
주제로 한 채팅이 시작된다.

Google Maps SDK\,[7]를 통해 다음 방문 장소까지의 내비게이션을 제공하며,
지도 화면에는 day별 이동 경로 점선이 표시된다. 장소 상세 정보는
\texttt{GET /tourist/info}를 통해 KTO 한국관광공사 TourAPI\,[10]에서 사진·정보를
7일 캐시 기준으로 조회한다.

\subsection{여행 일지 자동 생성}

탐험 종료 후 방문 장소 요약 화면을 거쳐 여행 일지를 자동 생성할 수 있다.
방문 장소 이력, 도착 시각, 마을 할망과 나눈 대화에서 언급된 설화를 구조화된
컨텍스트로 GPT-4o에 제공한다. 생성 결과는 300\textasciitilde500자 한국어 산문으로,
방문 장소를 시간 순서로 엮고 설화를 자연스럽게 녹인 여행자 1인칭 회고
형식이다. 생성된 일지는 iOS ShareLink를 통해 외부 공유가 가능하다.

% ── IV. 설계 변경 이력 ──────────────────────────────────
\section{설계 변경 이력}

\subsection{플랫폼: Flutter → iOS Swift/SwiftUI}

기초 조사서에는 Flutter 크로스플랫폼으로 계획하였으나, 이전 산책 앱
개발 경험에서 Flutter의 위치 기반 기능 일부를 결국 Kotlin·Swift 네이티브로
별도 구현해야 했던 교훈을 반영하여 iOS 네이티브 개발로 전환하였다.
CoreLocation 지오펜스, 백그라운드 GPS, iOS 알림 딥링크 등 위치 기반 기능의
정확도와 통합성이 향상되었다.

\subsection{코스 추천 방식: 3단계 재설계}

코스 추천 로직은 세 차례의 재설계를 거쳤다.

\textbf{1차 (설화→코스 역설계).}\quad 설화 6노드 LangGraph로 설화 배경 장소를
먼저 추출하고 Visit Jeju로 동선을 보완하는 방식. 설화를 코스의 뼈대로
삼았으나, 설화 위치 분포가 희박할 때 코스가 비현실적으로 구성되는 문제가
있었다.

\textbf{2차 (설화 취향 점수 기반).}\quad 사용자 설화 취향을 내부 7개 카테고리
점수로 변환하여 코스를 필터링하는 구조. 카테고리 분류가 기획 관점에서 설계되어
실제 설화 데이터의 다양성을 충분히 반영하지 못하는 한계가 있었다.

\textbf{3차 (코스→설화 매핑, 현재).}\quad 사용자가 여행을 계획할 때 가장 먼저
정하는 지역·기간을 입력받아 Visit Jeju 실제 코스를 먼저 선별하고, 그 위에
설화를 레이어로 얹는 방식. 코스의 현실성은 실제 데이터가 보장하고, 설화는
코스를 풍부하게 만드는 역할로 분리하여 사용자 경험이 안정적으로 개선되었다.

% ── V. 구현 현황 ─────────────────────────────────────────
\section{구현 현황}

중간 보고 시점 기준으로 서비스의 핵심 기능이 모두 구현 완료되었다.
iOS 앱은 취향 입력·코스 선택·설화 미리보기로 이어지는 계획 플로우,
GPS 도착 감지와 경로 안내를 포함한 탐험 플로우, 일별 방문 요약과
AI 설화 동행자 채팅, 여행 일지 자동 생성, 앱 세션 복원 기능을 갖춘다.
백엔드는 코스 추천(2단계 파이프라인), AI 설화 동행자 채팅, 여행 일지
생성, 설화 핀 조회, 관광정보 조회, TTS 변환의 API가 운영 중이다.

% ── VI. 결론 ────────────────────────────────────────────
\section{결 론}

중간 보고 시점 기준으로, 제주 설화·민담 여행 코스 추천 서비스의 핵심 기능
파이프라인이 완성되었다. 코스 추천은 세 차례의 설계 재검토를 통해
사용자 중심의 안정적인 2단계 구조로 확정되었으며, 여행 중 경험은 AI 설화
동행자 ``마을 할망''을 통해 설화와 대화가 자연스럽게 연결된다. 여행 후에는
GPT-4o 기반 여행 일지로 경험이 기록으로 남는다.

초기 설계 대비 플랫폼과 추천 방식이 변경되었으나, 핵심 가치인
\textbf{``실제 여행 코스 위에 제주 설화를 얹어 무형 문화유산의 접근성을
높인다''} 는 방향은 유지된다. 향후 사용자 평가를 통해 설화 콘텐츠 전달의
실효성을 검증할 예정이다.

% ── 참고문헌 ────────────────────────────────────────────
\section*{참 고 문 헌}

\begin{small}
\begin{enumerate}
\setlength{\itemsep}{1pt}
\item P. Lewis et al., ``Retrieval-Augmented Generation for
  Knowledge-Intensive NLP Tasks,'' \textit{NeurIPS}, 2020.
\item 제주특별자치도, 제주 설화정보 공공데이터 API (E07), 2023.
\item 제주특별자치도, 제주 민담정보 공공데이터 API (E08), 2023.
\item 제주관광공사, 제주관광정보시스템(VISIT JEJU) 여행장소 및
  여행세부일정 공공데이터, 2026.
\item Chroma, ``Chroma: the AI-native open-source embedding database,''
  2023. [Online]. Available: \textit{https://www.trychroma.com}
\item OpenAI, ``Text Embedding Models,''
  \textit{platform.openai.com/docs/guides/embeddings}, 2024.
\item Google, ``Google Maps Platform -- Geocoding API \& Directions
  API,'' \textit{developers.google.com/maps}, 2024.
\item S. Yao et al., ``ReAct: Synergizing Reasoning and Acting in
  Language Models,'' \textit{ICLR}, 2023.
\item LangChain, ``LangGraph: Build Stateful Multi-Actor
  Applications,'' \textit{langchain-ai.github.io/langgraph}, 2024.
\item 한국관광공사(KTO), TourAPI 공공데이터 서비스, 2024.
\end{enumerate}
\end{small}

\vspace{6pt}
\noindent\rule{\columnwidth}{0.4pt}
\vspace{4pt}
{\footnotesize $^\dagger$ 본 연구는 경희대학교 컴퓨터공학과 졸업프로젝트로 수행되었음.}

\end{document}
